The important case of \emph{binary} error correcting codes is still
open. In the binary case, the Gilbert-Varshamov bound gives the best
known (explicit or non-explicit) codes to date. For codes with
distance close to half, the bound shows that random linear codes
of length $n=O(k/\eps^2)$ are $[n,k,(1/2-\eps)n]_2$ codes. Finding
an explicit construction that attains this bound is an open problem.

Another closely related question is that of finding an
$[n,k,(1/2-\eps)n]_2$ binary code, in which the relative weight of
every non-zero codeword is in the range $[1/2-\eps,1/2+\eps]$.
Such codes are called \emph{$\eps$-balanced} and they are related to
another kind of combinatorial objects called \emph{$\eps$-biased
sets}. An $\eps$-biased set is a set $S \subseteq \set{0,1}^k$ such
that for every non-empty subset $T \subseteq [k]$, the binary random
variable $\bigoplus_{i \in T} s_i$, where $s$ is sampled uniformly
from $S$, has bias at most $\eps$. It turns out that $\eps$-biased
sets are just $\eps$-balanced codes in a different guise: the
%
%
rows of a matrix whose columns generate an $\eps$-balanced code form an $\eps$-biased set, and vice
versa. In terms of parameters, an $[n,k]_2$ $\eps$-balanced code is
equivalent to an $\eps$-biased set $S \subseteq \set{0,1}^k$ of size
$n$.

The status of $\eps$-balanced codes is similar to that of
$[n,k,(1/2-\eps)n]_2$ codes. In both cases the probabilistic
method gives non-explicit $[n,k]_2$ $\eps$-balanced codes with
$n=O({k / \eps^2})$, whereas the best lower bound is
\[
n=\Omega\left(\frac{k}{\eps^2 \logeps}\right)\,.
\] For a discussion of these
bounds see~\cite[Section 7]{AGHP:92}.
